\documentclass[10pt,a4paper]{amsart}
\usepackage[margin=19mm]{geometry}
\usepackage{amsmath,amssymb,mathtools}
\usepackage[scr=rsfs,cal=euler]{mathalfa}
\usepackage{xfrac}
\usepackage[unicode,hidelinks,bookmarks=false]{hyperref}
\newcommand{\ie}{i.e.\ }
\newcommand{\cf}{cf.\ }
\newcommand{\resp}{respectively\ }
\newcommand{\Subsection}{\subsection}
\providecommand{\nobreakdash}{\nobreak\hbox{-}}
\setlength{\emergencystretch}{2em}
\multlinegap0pt
\usepackage{dsfont}
\usepackage{amssymb}
\usepackage{mathtools}
\usepackage{tikz}
\newtheorem{lemma}{Lemma}
\newtheorem{definition}{Definition}
\newtheorem{proposition}{Proposition}
\newcommand{\CF}{\mathcal F}
\newcommand{\E}{\mathbb E}
\newcommand{\F}{{\mathcal F}}
\title{Kohn--Nirenberg quantization of the affine group and related examples}
\author{Pierre Bieliavsky, Victor Gayral, Sergey Neshveyev, Lars Tuset}
\date{Source: arXiv:2604.08274v1 (CC BY 4.0)}
\begin{document}
\maketitle

\noindent\emph{Source and licence.} Kohn--Nirenberg quantization of the affine group and related examples. Version arXiv:2604.08274v1 (CC BY 4.0), distributed by arXiv under the Creative Commons Attribution 4.0 International licence, http://creativecommons.org/licenses/by/4.0/, as recorded in the arXiv licence metadata for this exact version and shown on the abstract page https://arxiv.org/abs/2604.08274v1. Full licence notice: \texttt{licenses/arxiv-2604.08274-CC-BY-4.0.txt}.\par\smallskip

\noindent\textit{Editorial scope.} Counted unit: the article's proposition CRF (source0.tex lines 1348--1354). The article's complete original proof of it is delivered with it (source0.tex lines 1355--1450, one \texttt{proof} environment of 96 lines). No mathematics is written, repaired or re-derived: the statement and the proof are the article's own lines, in the article's own notation, and no retained line is substituted or re-flowed.  Retained uncounted as the source-local context the proof needs: lines 984--987; lines 1245--1252; lines 1316--1320; lines 1337--1345.  Nothing was omitted from any retained span, so the package carries no omission record.  No retained line contains a citation, so the wrapper carries no bibliography and no bibliography entry.  No retained line contains a figure environment, an image-inclusion command or drawing code, so no image is delivered and the package declares no asset dependency.  Editorial formatting only, in addition: an AMS \texttt{amsart} preamble that restates the article's own declarations and macros used by the retained lines, namely the environments \texttt{lemma}, \texttt{definition}, \texttt{proposition} on separate counters, together with the standard packages those lines need.  The retained environments print in this document as Lemma 1, Definition 1, Proposition 1 in order of appearance, whereas the article prints the same items with the numbering of its own full document.  The complete native source of the article as distributed in the arXiv e-print for this exact version is bundled unchanged as \texttt{sources/198148/source0.tex}.
 \begin{equation}
 \label{phi}
\phi_\ell:Q_\ell\to \hat V_\ell, \quad q_\ell\mapsto q_\ell^\flat\xi_{0,\ell},
\end{equation}

\begin{lemma}
\label{RecE}
We have the following inductive relation:
$$
\E_\ell(p_\ell ,n_\ell)=e^{i\langle \phi_\ell(q_\ell),v_\ell\rangle}\,\E_{\ell-1}\big(p_{\ell-1} ,\tilde\alpha_{q_\ell^{-1}}(n_{\ell-1})\big),
$$
where the map $\phi_\ell:Q_\ell\to\hat V_\ell$ is given by \eqref{phi}.
\end{lemma}

\begin{align}
\label{TRL}
\tilde\alpha_{ p_\ell^{-1}}(n_\ell)
={\bf C}_{ p_\ell^{-1}}(v_\ell)\tilde\alpha_{ p_{\ell-1}^{-1}}(\tilde\alpha_{q_\ell^{-1}}(n_{\ell-1})),
\end{align}

\begin{definition}
\label{URPN}
Let  $U_{\tilde \alpha_\ell}$ be the unitary representation of $P_\ell$   on $L^2(N_\ell)$ and let $U_{\tilde\beta_\ell}$ be the unitary representation of  $N_\ell$
 on $L^2(P_\ell)$ given by
$$
U_{\tilde \alpha_\ell}( p_\ell)f(n_\ell):=J_{\tilde\alpha_\ell}(p_\ell^{-1},n_\ell)^{1/2}\, f\big(\tilde\alpha_{p_\ell^{-1}}(n_\ell)\big)
\quad\mbox{and}\quad U_{\tilde \beta_\ell}( n_\ell)f(p_\ell):= f\big(\tilde\beta_{n_\ell^{-1}}(p_\ell)\big).
$$
\end{definition}

\begin{proposition}
\label{CRF}
For $(p_\ell,n_\ell)\in P_\ell\times N_\ell$, we have
$$
\lambda_{  p_\ell}\,\F_\ell=\F_\ell\, U_{\tilde \alpha_\ell}(p_\ell) \quad\mbox{and}\quad \F_\ell\, \lambda_{n_\ell}=\overline\E_\ell(\cdot,n_\ell)\,U_{\tilde \beta_\ell}(n_\ell)\,\F_\ell.
$$
\end{proposition}

\begin{proof}
In addition to the representations appearing in Definition \ref{URPN}, let us consider the unitary representations $C_{V_\ell}$ and $C_{\hat V_\ell}$ of $H_\ell$ on
$L^2(V_\ell)$ and $L^2(\hat V_\ell)$, respectively, and the unitary representation $C_{Q_\ell}$ of~$P_{\ell-1}$ on~$L^2 (Q_\ell )$ defined by
\begin{align*}
C_{V_\ell}(h_\ell)f(v_\ell)&:=|h_\ell|_{V_\ell}^{-1/2}\,f({\bf C}_{h_\ell^{-1}}(v_\ell)),\\
C_{\hat V_\ell}(h_\ell)f(\xi_\ell)&:=|h_\ell|_{V_\ell}^{1/2}\,f({h_\ell^{-1}}^\flat\xi_\ell),\\
C_{Q_\ell}(p_{\ell-1})f(q_\ell)&:=|p_{\ell-1}|_{Q_\ell}^{-1/2}\,f({\bf C}_{p_{\ell-1}^{-1}}(q_\ell)).
\end{align*}
Lastly, we let $U_{\tilde\alpha,N_{\ell-1}}$ and $U_{\tilde\beta,Q_\ell}$ be the unitary representations of $Q_\ell$  on $L^2(N_{\ell-1})$  and  of $N_{\ell-1}$  on $L^2(Q_\ell )$ defined by
\begin{align*}
U_{\tilde\alpha,N_{\ell-1}}( q_\ell)f(n_{\ell-1})&:=|q_\ell|_{V_\ell}^{1/2}\,J_{\tilde\alpha_\ell}(q_\ell^{-1},n_{\ell-1})^{1/2}\, f\big(\tilde\alpha_{q_\ell^{-1}}(n_{\ell-1})\big),\\
U_{\tilde\beta,Q_\ell}(n_{\ell-1} )f(q_\ell)&:= f\big(\tilde\beta_{n_{\ell-1}^{-1}}(q_\ell)\big).
\end{align*}

Since the map $\phi_\ell:Q_\ell\to\hat V_\ell$ intertwines the left action of $Q_\ell$ on itself with the dual action on~$\hat V_\ell$,
 we deduce
 \begin{align}
 \label{CR1}
  \lambda_{q_\ell}\, U_{\phi_\ell}\,\F_{V_\ell}= U_{\phi_\ell}\,C_{\hat V_\ell}(q_\ell)\F_{V_\ell}= U_{\phi_\ell}\,\F_{V_\ell}\, C_{V_\ell}(q_\ell).
  \end{align}
This already proves the first relation for   $\ell=1$,  since in this case we have $U_{\tilde \alpha_1}=C_{V_1}$ and  $\CF_1=V_{\phi_1}\,\F_{V_1}$.

To prove the first commutation relation for all $\ell$, we proceed by induction and first consider the case where $p_\ell$ belongs to $Q_\ell$. From
the identities $ \lambda_{  q_\ell}V_{1,\ell}^*=V_{1,\ell}^*( \lambda_{q_\ell}\otimes 1)$, $( \lambda_{q_\ell}\otimes 1)  V_{\tilde\alpha_\ell} = V_{\tilde\alpha_\ell} (
  \lambda_{q_\ell}\otimes U_{\tilde\alpha,N_{\ell-1}}(q_\ell))$ and \eqref{CR1}, we get
 $$
 \lambda_{  q_\ell}\,\F_\ell=V_{1,\ell}^*\,(1\otimes \F_{\ell-1})\, V_{\tilde\alpha_\ell} \,
 (  U_{\phi_\ell}\,\F_{V_\ell}\otimes 1)(C_{V_\ell}(q_\ell)\otimes U_{\tilde\alpha,N_{\ell-1}}(q_\ell))\,V_{2,\ell},
$$
 and this is what we need,  because we have
$ \tilde\alpha_{ p_\ell^{-1}}(n_\ell)
={\bf C}_{ q_\ell^{-1}}(v_\ell)\tilde\alpha_{q_\ell^{-1}}(n_{\ell-1})$ by \eqref{TRL}.

Next, consider the case where  $p_\ell$ belongs to $P_{\ell-1}$.  We have
$\lambda_{p_{\ell-1}}V_{1,\ell}^*=V_{1,\ell}^*(C_{Q_\ell}(p_{\ell-1})\otimes \lambda_{p_{\ell-1}})$,
and thus we get by the induction hypothesis:
$$
\lambda_{p_{\ell-1}}\,\F_\ell=V_{1,\ell}^*\,(1\otimes \F_{\ell-1})\, (C_{Q_\ell}(p_{\ell-1})\otimes U_{\tilde\alpha_{\ell-1}}(p_{\ell-1}))\,
 V_{\tilde\alpha_\ell} \,( U_{\phi_\ell}\,\F_{V_\ell}\otimes 1)\,V_{2,\ell}.
$$
It is easy to see that  $ C_{Q_\ell} (p_{\ell-1})\otimes U_{\tilde\alpha_{\ell-1}}(p_{\ell-1})$ commutes with
$ V_{\tilde\alpha_\ell} $, and since ${p_{\ell-1}^{-1}}^\flat\xi_{0,\ell}=\xi_{0,\ell}$, we have
$ C_{Q_\ell} (p_{\ell-1})\, U_{\phi_\ell}\,\F_{V_\ell}= U_{\phi_\ell}\, C_{\hat V_\ell} (p_{\ell-1})\,\F_{V_\ell}= U_{\phi_\ell}\,\F_{V_\ell}\, C_{V_\ell}(p_{\ell-1})$. Hence we obtain:
 $$
\lambda_{p_{\ell-1}}\,\F_\ell=V_{1,\ell}^*\,(1\otimes \F_{\ell-1})\,
 V_{\tilde\alpha_\ell} \,( U_{\phi_\ell}\,\F_{V_\ell}\otimes 1)( C_{Q_\ell} (p_{\ell-1})\otimes U_{\tilde\alpha_{\ell-1}}(p_{\ell-1}))\,V_{2,\ell},
$$
and we conclude again by \eqref{TRL}.

\smallskip

 Let us now prove the second commutation relation. When $\ell=1$, the result follows  from
$U_{\tilde \beta_1}(v_1)={\rm Id}$ and $\E_1(q_1,v_1)=e^{i\langle\phi_1(q_1),v_1\rangle}$. We then proceed by induction on $\ell$ and first consider the case where
 $n_\ell$ belongs to $V_\ell$. From the identities $V_{2,\ell}\,\lambda_{v_\ell}=(\lambda_{v_\ell}\otimes 1)\,V_{2,\ell}$,
$ U_{\phi_\ell}\,\F_{V_\ell}\,\lambda_{v_\ell}= e^{-i\langle\phi_\ell(\cdot),v_\ell\rangle}\, U_{\phi_\ell}\,\F_{V_\ell}$
and $ V_{\tilde\alpha_\ell} \,(e^{-i\langle\phi_\ell(\cdot),v_\ell\rangle}\otimes 1)=(e^{-i\langle\phi_\ell(\cdot),v_\ell\rangle}\otimes 1)\,  V_{\tilde\alpha_\ell} $,
we deduce that
$$
\F_\ell\,\lambda_{v_\ell}=V_{1,\ell}^*\,(e^{i\langle\phi_\ell(\cdot),v_\ell\rangle}\otimes 1)\,(1\otimes \F_{\ell-1})\, V_{\tilde\alpha_\ell} \,( U_{\phi_\ell}\,\F_{V_\ell}\otimes 1)\,V_{2,\ell},
$$
and the result follows, because  $\E_\ell(p_\ell, v_\ell)=e^{i\langle\phi_\ell(q_\ell),v_\ell\rangle}$.

Next, we consider the case where  $n_\ell$ belongs to $N_{\ell-1}$.
The relations $V_{2,\ell}\,\lambda_{n_{\ell-1}}=( C_{V_\ell} (n_{\ell-1})\otimes \lambda_{n_{\ell-1}})\,V_{2,\ell}$ and $ U_{\phi_\ell}\,\F_{V_\ell}\, C_{V_\ell} (n_{\ell-1})=
 U_{\phi_\ell}\, C_{\hat V_\ell} (n_{\ell-1})\,\F_{V_\ell}=U_{\tilde\beta,Q_\ell}(n_{\ell-1})\, U_{\phi_\ell}\,\F_{V_\ell}$ give
$$
\F_\ell\,\lambda_{n_{\ell-1}}=V_{1,\ell}^*\,(1\otimes \F_{\ell-1})\, V_{\tilde\alpha_\ell} \,(U_{\tilde\beta,Q_\ell}(n_{\ell-1})
\otimes \lambda_{n_{\ell-1}})\,( U_{\phi_\ell}\,\F_{V_\ell}\otimes 1)\,V_{2,\ell}.
$$
Observe now that
$$
 V_{\tilde\alpha_\ell} \,(U_{\tilde\beta,Q_\ell}( n_{\ell-1})
\otimes \lambda_{ n_{\ell-1}})f(\tilde q_\ell,\tilde n_{\ell-1})= |q_\ell|_{V_\ell}^{-1/2}\,
J_{\tilde\gamma_\ell}(\tilde q_\ell,n_{\ell-1})^{1/2}\,f\big(\tilde\beta_{n_{\ell-1}^{-1}}(\tilde q_\ell),n_{\ell-1}^{-1}\tilde\alpha_{\tilde q_\ell}(\tilde n_{\ell-1}) \big).
$$
Since
$$
n_{\ell-1}^{-1}\tilde\alpha_{\tilde q_\ell}(\tilde n_{\ell-1})=\tilde\alpha_{\tilde\beta_{n_{\ell-1}^{-1}}(\tilde q_\ell)}\big(\tilde\alpha_{\tilde q_\ell^{-1}}(n_{\ell-1})^{-1}\tilde n_{\ell-1}\big),
$$
we deduce that
$$
 V_{\tilde\alpha_\ell} \,(U_{\tilde\beta,Q_\ell}( n_{\ell-1})
\otimes \lambda_{ n_{\ell-1}})f(\tilde q_\ell,\tilde n_{\ell-1})= V_{\tilde\alpha_\ell} f\big(\tilde\beta_{n_{\ell-1}^{-1}}(\tilde q_\ell),\tilde\alpha_{\tilde q_\ell^{-1}}(n_{\ell-1})^{-1}\tilde n_{\ell-1} \big).
$$
Therefore we get by the induction hypothesis that
\begin{multline*}
(1\otimes \F_{\ell-1})\, V_{\tilde\alpha_\ell} \,(U_{\tilde\beta,Q_\ell}( n_{\ell-1})\otimes \lambda_{ n_{\ell-1}})f(q_\ell,p_{\ell-1})=\\
\overline\E_{\ell-1}\big(p_{\ell-1}, \tilde\alpha_{q_\ell^{-1}}(n_{\ell-1})\big)\,
(1\otimes \F_{\ell-1})\, V_{\tilde\alpha_\ell} f\big(\tilde\beta_{n_{\ell-1}^{-1}}(q_\ell),\tilde\beta_{\tilde\alpha_{ q_\ell^{-1}}(n_{\ell-1})^{-1} }(p_{\ell-1})\big).
\end{multline*}
A simple computation shows that
$$
U_{\tilde \beta_\ell}(n_{\ell-1})\,V_{1,\ell}^*f(p_\ell q_{\ell-1})=f\big(\tilde\beta_{n_{\ell-1}^{-1}}(q_\ell),\tilde\beta_{\tilde\alpha_{ q_\ell^{-1}}(n_{\ell-1})^{-1} }(p_{\ell-1})\big),
$$
and the conclusion follows from the equality $\E_{\ell-1}\big(p_{\ell-1}, \tilde\alpha_{q_\ell^{-1}}(n_{\ell-1})\big)=\E_\ell(p_\ell,n_{\ell-1})$ proven in Lemma \ref{RecE}.
\end{proof}

\end{document}
